% BOOK-FILE-ID: f-49d89f9a5544
\subsection{発展・特殊型}\label{節：発展・特殊型}
\sectionref{節：発展的な解法}と\sectionref{節：式の構造を使う}の15題である.
\bookterm{floor}{ガウス記号}には整数性・余り,多項間の関係には\bookterm{difference}{階差}・\bookterm{characteristic}{特性方程式},
\bookterm{nonlinear}{非線形}な式には対称式・\bookterm{trigonometric}{三角関数}の恒等式を使う.
候補は\bookterm{initial}{初期条件}と漸化式で確かめ,分母や角度の定義域も確認する.
問14・問15では初項を動かし,定\bookterm{sequence}{数列}・周期列・有界な\bookterm{sequence}{数列}が現れる条件を調べる.有界や$+\infty$への発散の意味は\sectionref*{節：数列の極限}を参照する.
以下,特に指定しない限り,$n$は正の整数とする.
$\lfloor x\rfloor$は$x$を超えない最大の整数を表す.
周期の最小値とは,すべての添字で$a_{n+p}=a_n$となる正の整数$p$の最小値をいう.

\subsubsection{問題}\label{節：発展・特殊型：問題}
\begin{enumerate}[label=\textbf{問\arabic*.},leftmargin=3.8em,format=\bookitemlink{practice-special}{problem}{answer}]
  % E1: ガウス型・整数の半減
  \item \bookterm{sequence}{数列}$\{a_n\}$を
    \[
      a_1=2026,\qquad
      a_{n+1}=\left\lfloor\frac{a_n}{2}\right\rfloor
    \]
    によって定める.
    \par\smallskip\noindent
    (1) \bookterm{general}{一般項}を求めよ.
    \par\smallskip\noindent
    (2) 初めて$a_n=0$となる添字$n$を求めよ.
    \par\smallskip\noindent
    (3) 非負整数$M$に対し,別の数列を
    \[
      b_1=M,\qquad
      b_{n+1}=\left\lfloor\frac{b_n+6}{2}\right\rfloor
    \]
    と定める.$b_n$の一般項と,最終的に一定になる値を求めよ.
    % E2: ガウス型・三つずつの群
  \item \bookterm{sequence}{数列}$\{a_n\}$を
    \[
      a_0=0,\qquad
      a_{n+1}=a_n+\left\lfloor\frac{n+1}{3}\right\rfloor
      \quad(n\ge0)
    \]
    によって定める.\,
    (1) $n=3q+r$（$q\ge0$,$r=0,1,2$）とおき,
    \bookterm{general}{一般項}を$q,r$で表せ.
    \par\smallskip\noindent
    (2) $a_n$を,次の範囲にある格子点（座標がともに整数の点）の個数として説明せよ.
    \[
      1\le x\le n,\qquad 1\le y\le\frac{x}{3}.
    \]
    % E3: ガウス型・余りの有限個の状態
  \item \bookterm{sequence}{数列}$\{a_n\}$を
    \[
      a_1=1,\qquad
      a_{n+1}=2a_n+1
      -5\left\lfloor\frac{2a_n+1}{5}\right\rfloor
    \]
    によって定める.
    \par\smallskip\noindent
    (1) すべての項が$0,1,2,3,4$のいずれかになることを示せ.
    \par\smallskip\noindent
    (2) \bookterm{general}{一般項}を,$n$を$4$で割った余りで場合分けして表せ.
    また,周期の最小値を求めよ.
    % E4: 隣接4項間・三階差が定数
  \item \bookterm{sequence}{数列}$\{a_n\}$を
    \[
      a_1=2,\qquad a_2=9,\qquad a_3=28,
    \]
    \[
      a_{n+3}=3a_{n+2}-3a_{n+1}+a_n+6
    \]
    によって定める.
    \par\smallskip\noindent
    (1) $d_n=a_{n+2}-2a_{n+1}+a_n$とおき,$d_n$を求めよ.
    \par\smallskip\noindent
    (2) \bookterm{general}{一般項}を求めよ.
    % E5: 隣接4項間・異なる三つの特性根
  \item \bookterm{sequence}{数列}$\{a_n\}$を
    \[
      a_1=3,\qquad a_2=6,\qquad a_3=14,
    \]
    \[
      a_{n+3}=6a_{n+2}-11a_{n+1}+6a_n
    \]
    によって定める.
    \par\smallskip\noindent
    (1) $d_n=a_{n+2}-3a_{n+1}+2a_n$とおき,$d_n$を求めよ.
    \par\smallskip\noindent
    (2) \bookterm{general}{一般項}を求めよ.
    % E6: 虚数解型・周期6
  \item \bookterm{sequence}{数列}$\{a_n\}$を
    \[
      a_1=0,\qquad a_2=1,\qquad
      a_{n+2}=a_{n+1}-a_n
    \]
    によって定める.
    \par\smallskip\noindent
    (1) \bookterm{trigonometric}{三角関数}を用いて\bookterm{general}{一般項}を表せ.
    \par\smallskip\noindent
    (2) 周期の最小値を求めよ.
    % E7: 虚数解型・振幅が減衰する数列
  \item \bookterm{sequence}{数列}$\{a_n\}$を
    \[
      a_1=1,\qquad a_2=0,\qquad
      a_{n+2}=a_{n+1}-\frac12a_n
    \]
    によって定める.\bookterm{trigonometric}{三角関数}を用いて\bookterm{general}{一般項}を表し,
    $\lim_{n\to\infty}a_n$を求めよ.
    % E8: 非隣接2項間・奇数列と偶数列の逆数
  \item \bookterm{sequence}{数列}$\{a_n\}$を
    \[
      a_1=1,\qquad a_2=2,\qquad
      a_{n+2}=\frac{a_n}{1+a_n}
    \]
    によって定める.
    \par\smallskip\noindent
    (1) すべての項が正であることを示し,
    $a_{2k-1},a_{2k}$を$k$で表せ.
    ただし,$k$は正の整数とする.
    \par\smallskip\noindent
    (2) $\lim_{n\to\infty}a_n$を求めよ.
    % E9: 基本対称式型・平行移動して恒等式を使う
  \item \bookterm{sequence}{数列}$\{a_n\}$を
    \[
      a_1=\frac72,\qquad
      a_{n+1}=a_n^2-2a_n
    \]
    によって定める.
    \par\smallskip\noindent
    (1) $b_n=a_n-1$とおき,$b_n$の\bookterm{recurrence}{漸化式}を求めよ.
    \par\smallskip\noindent
    (2) 恒等式
    \[
      \left(u+u^{-1}\right)^2-2=u^2+u^{-2}
      \quad(u\ne0)
    \]
    を用いて$a_n$の\bookterm{general}{一般項}を求めよ.
    % E10: 基本対称式型・三乗の対称式
  \item \bookterm{sequence}{数列}$\{a_n\}$を
    \[
      a_1=\frac52,\qquad
      a_{n+1}=a_n^3-3a_n
    \]
    によって定める.
    \par\smallskip\noindent
    (1) 恒等式
    \[
      \left(u+u^{-1}\right)^3
      -3\left(u+u^{-1}\right)=u^3+u^{-3}
    \]
    を示せ.ただし,$u\ne0$とする.
    \par\smallskip\noindent
    (2) \bookterm{general}{一般項}を求めよ.
    % E11: 三角関数型・余弦の倍角と周期3
  \item \bookterm{sequence}{数列}$\{a_n\}$を
    \[
      a_1=\cos\frac{2\pi}{7},\qquad
      a_{n+1}=2a_n^2-1
    \]
    によって定める.
    \par\smallskip\noindent
    (1) \bookterm{general}{一般項}を求めよ.
    \par\smallskip\noindent
    (2) 周期の最小値を求めよ.
    % E12: 三角関数型・正接の定義域と周期4
  \item 初項と次の項を計算する規則を
    \[
      a_1=\tan\frac{\pi}{5},\qquad
      a_{n+1}=\frac{2a_n}{1-a_n^2}
    \]
    とする.
    \par\smallskip\noindent
    (1) \bookterm{general}{一般項}の候補を求め,すべての段階で分母が$0$にならず,
    \bookterm{sequence}{数列}が定義できることを示せ.
    \par\smallskip\noindent
    (2) 周期の最小値を求めよ.
    % E13: ニュートン法型・初項が収束先より小さい場合
  \item \bookterm{sequence}{数列}$\{a_n\}$を
    \[
      a_1=1,\qquad
      a_{n+1}=\frac{a_n^2+3}{2a_n}
    \]
    によって定める.
    \par\smallskip\noindent
    (1) すべての項が正であることを示せ.
    \par\smallskip\noindent
    (2) $b_n=\frac{a_n-\sqrt3}{a_n+\sqrt3}$とおき,
    $b_n,a_n$の\bookterm{general}{一般項}を求めよ.
    \par\smallskip\noindent
    (3) $n\ge2$で$a_n>\sqrt3$かつ
    $a_{n+1}<a_n$であることを示し,$\lim_{n\to\infty}a_n$を求めよ.
    % E14: 初項依存型・不動点と2周期点
  \item 実数$c$を初項として,\bookterm{sequence}{数列}$\{a_n\}$を
    \[
      a_1=c,\qquad a_{n+1}=a_n^2-1
    \]
    によって定める.
    \par\smallskip\noindent
    (1) \bookterm{sequence}{数列}が定\bookterm{sequence}{数列}となる$c$をすべて求めよ.
    \par\smallskip\noindent
    \BookMathLayout{(2) 定\bookterm{sequence}{数列}ではなく,すべての$n$で$a_{n+2}=a_n$を満たす
      $c$をすべて求め,それぞれの場合の\bookterm{general}{一般項}を表せ.}{(2) 定\bookterm{sequence}{数列}ではなく,\\すべての$n$で$a_{n+2}=a_n$を満たす
      $c$をすべて求め,それぞれの場合の\bookterm{general}{一般項}を表せ.}
    % E15: 初項依存型・有界になる初項の範囲
  \item 実数$c$を初項として,\bookterm{sequence}{数列}$\{a_n\}$を
    \[
      a_1=c,\qquad a_{n+1}=a_n^2-2
    \]
    によって定める.
    \par\smallskip\noindent
    (1) $c=-2,-1,0,2$のそれぞれについて,最初の4項を求めよ.
    \par\smallskip\noindent
    (2) $-2\le c\le2$のとき,$c=2\cos\theta$
    （$0\le\theta\le\pi$）とおき,\bookterm{general}{一般項}を求めよ.
    また,すべての項が$-2$以上$2$以下であることを示せ.
    \par\smallskip\noindent
    (3) $|c|>2$のとき,$a_n\to+\infty$を示せ.
    以上から,\bookterm{sequence}{数列}が有界となる$c$の範囲を求めよ.
\end{enumerate}

\subsubsection{方針}\label{節：発展・特殊型：方針}
\begin{enumerate}[label=\textbf{問\arabic*.},leftmargin=3.8em,format=\bookitemlink{practice-special}{hint}{problem}]
  % E1
  \item 整数を$2$で割った商と余りに分け,
    $\left\lfloor\frac{\lfloor x\rfloor}{2}\right\rfloor
      =\left\lfloor\frac{x}{2}\right\rfloor$となることを確かめる.
    初めて$0$になる位置は,$2026$と$2$の累乗を比べればよい.
    (3)では$b_n-6$を考え,負の数の切り捨てにも注意する.
    % E2
  \item \bookterm{difference}{階差}を足すと$\sum_{k=1}^{n}\lfloor\frac{k}{3}\rfloor$になる.
    同じ値を取る添字を三つずつまとめ,最後の群を$r+1$項として数える.
    $\lfloor\frac{k}{3}\rfloor$は,$x=k$上にある指定範囲の格子点の個数でもある.
    % E3
  \item $x-5\lfloor\frac{x}{5}\rfloor$は整数$x$を$5$で割った余りである.
    項の値が五つに絞れたら,それぞれが次にどの値へ移るかを表にする.
    % E4
  \item $d_{n+1}-d_n$を作ると元の\bookterm{recurrence}{漸化式}を移項した式の左辺が現れる.
    その後,$e_n=a_{n+1}-a_n$を求め,もう一度足して$a_n$へ戻す.
    % E5
  \item $d_{n+1}-3d_n$を展開すると$0$になる.
    得られた3項間の式を満たす$3^{n-1}$を引き,残りを解く.
    % E6
  \item \bookterm{characteristic}{特性方程式}の根の絶対値と\bookterm{argument}{偏角}を調べる.
    最初の二つの項に合う$\sin$の式を作り,実際の値の並びから周期を確かめる.
    % E7
  \item 特性方程式の解の絶対値は$1$未満である.
    \bookterm{polar}{極形式}から\bookterm{trigonometric}{三角関数}の\bookterm{general}{一般項}を作り,その絶対値を\bookterm{geometric}{等比数列}で上から押さえる.
    % E8
  \item 正値を確かめてから逆数を取る.
    奇数番目と偶数番目を別々の\bookterm{sequence}{数列}にすると,二つの\bookterm{arithmetic}{等差数列}が得られる.
    % E9
  \item $a_n^2-2a_n=(a_n-1)^2-1$と\bookterm{square}{平方完成}する.
    $b_1=2+2^{-1}$と書けば,恒等式によって指数が毎回2倍になる.
    % E10
  \item 展開すると,交差項が$3(u+u^{-1})$であることが分かる.
    初項を$2+2^{-1}$と表し,指数を毎回3倍にする候補を作る.
    % E11
  \item 余弦の倍角公式で,角度を毎回2倍にする候補を作る.
    3回倍にした角度と元の角度との差が$2\pi$の整数倍になることを調べる.
    % E12
  \item 正接の倍角公式を使う.
    候補の角度が$\frac{\pi}{2}$に$\pi$の整数倍を加えた値になるかを,
    整数の偶奇で判断する.次の角度でも正接が定義できれば,\bookterm{recurrence}{漸化式}の分母も$0$にならない.
    % E13
  \item $a_{n+1}\pm\sqrt3$の分子は平方になる.比を取って$b_{n+1}=b_n^2$を導く.
    $b_1<0$なので,\bookterm{logarithm}{対数}を取らずにべき乗の形を帰納法で確かめる.
    減少を示す範囲が$n\ge2$であることにも注意する.
    % E14
  \item 定\bookterm{sequence}{数列}には$a_2=a_1$,周期2の候補には$a_3=a_1$が必要である.
    後者の方程式を因数分解し,定\bookterm{sequence}{数列}の初項を除く.残った初項が実際に二つの値を往復することも確かめる.
    % E15
  \item $|c|\le2$では$2\cos\theta$という表示を使う.
    $|c|>2$では第2項が$2$を超えるので,
    $a_{n+1}-2=(a_n-2)(a_n+2)$から下からの評価を作る.
\end{enumerate}

\subsubsection{答え}\label{節：発展・特殊型：答え}
\begin{enumerate}[label=\textbf{問\arabic*.},leftmargin=3.8em,format=\bookitemlink{practice-special}{answer}{problem}]
  % E1
  \item (1)
    \[
      a_n=\left\lfloor\frac{2026}{2^{n-1}}\right\rfloor.
    \]
    (2) $n=12$である.
    \par\smallskip\noindent
    (3)
    \[
      b_n=6+\left\lfloor\frac{M-6}{2^{n-1}}\right\rfloor.
    \]
    有限回の操作の後で,$M<6$なら$5$,$M\ge6$なら$6$に固定される.
    % E2
  \item (1)
    \[
      a_n=\frac{3q(q-1)}2+q(r+1),
    \]
    \[
      n=3q+r,\qquad q\ge0,\quad r=0,1,2.
    \]
    (2) $x=k$ごとに$\lfloor\frac{k}{3}\rfloor$個の点があり,その和が$a_n$である.
    % E3
  \item (1) 整数性と\bookterm{floor}{ガウス記号}の定義から,各項は$0$以上$4$以下の整数となる.
    \par\smallskip\noindent
    (2) $k$を$0$以上の整数として,
    \[
      a_{4k+1}=1,\quad a_{4k+2}=3,
    \]
    \[
      a_{4k+3}=2,\quad a_{4k+4}=0.
    \]
    周期の最小値は$4$である.
    % E4
  \item (1) $d_n=6n+6$.
    \par\smallskip\noindent
    (2) $a_n=n^3+1$.
    % E5
  \item (1) $d_n=2\cdot3^{n-1}$.
    \par\smallskip\noindent
    (2)
    \[
      a_n=1+2^{n-1}+3^{n-1}.
    \]
    % E6
  \item (1)
    \[
      a_n=\frac2{\sqrt3}\sin\frac{(n-1)\pi}{3}.
    \]
    (2) 周期の最小値は$6$である.
    % E7
  \item
    \[
      a_n=\left(\frac1{\sqrt2}\right)^{n-1}
      \left(\cos t_n-\sin t_n\right),
    \]
    \[
      t_n=\frac{(n-1)\pi}{4},\qquad
      \lim_{n\to\infty}a_n=0.
    \]
    % E8
  \item (1) すべての項は正であり,
    \[
      a_{2k-1}=\frac1k,\qquad
      a_{2k}=\frac2{2k-1}\quad(k\ge1).
    \]
    (2) \bookterm{limit}{極限}は$0$である.
    % E9
  \item (1) $b_1=\frac52$,$b_{n+1}=b_n^2-2$.
    \par\smallskip\noindent
    (2)
    \[
      a_n=1+2^{\,2^{n-1}}+2^{-2^{n-1}}.
    \]
    % E10
  \item (1) 左辺を展開すると右辺を得る.
    \par\smallskip\noindent
    (2)
    \[
      a_n=2^{\,3^{n-1}}+2^{-3^{n-1}}.
    \]
    % E11
  \item (1)
    \[
      a_n=\cos\left(2^{n-1}\frac{2\pi}{7}\right).
    \]
    (2) 周期の最小値は$3$である.
    % E12
  \item (1)
    \[
      a_n=\tan\left(2^{n-1}\frac{\pi}{5}\right).
    \]
    すべての項で正接が定義でき,\bookterm{recurrence}{漸化式}の分母も$0$にならない.
    \par\smallskip\noindent
    (2) 周期の最小値は$4$である.
    % E13
  \item (1) 初項が正で,正の項から次の正の項が定まる.
    \par\smallskip\noindent
    (2) $q=\sqrt3-2$とおくと,
    \[
      b_n=q^{\,2^{n-1}},\qquad
      a_n=\sqrt3\,\frac{1+q^{\,2^{n-1}}}{1-q^{\,2^{n-1}}}.
    \]
    (3) $n\ge2$で設問の二つの不等式が成り立ち,
    \[
      \lim_{n\to\infty}a_n=\sqrt3.
    \]
    % E14
  \item (1)
    \[
      c=\frac{1+\sqrt5}{2},\quad
      \frac{1-\sqrt5}{2}.
    \]
    それぞれ$a_n=c$である.
    \par\smallskip\noindent
    (2) $c=0,-1$であり,どちらも
    \[
      a_n=-\frac12+\left(c+\frac12\right)(-1)^{n-1}
    \]
    と表せる.周期の最小値は$2$である.
    % E15
  \item (1)
    \[
      \begin{array}{c|rrrr}
        c&a_1&a_2&a_3&a_4\\ \hline
        -2&-2&2&2&2\\
        -1&-1&-1&-1&-1\\
        0&0&-2&2&2\\
        2&2&2&2&2
      \end{array}
    \]
    (2)
    \[
      a_n=2\cos(2^{n-1}\theta),\qquad -2\le a_n\le2.
    \]
    (3) $n\ge2$で
    \[
      a_n\ge2+4^{n-2}(c^2-4)
    \]
    となり,$a_n\to+\infty$である.
    有界となる範囲は$-2\le c\le2$である.
\end{enumerate}

\subsubsection{解説}\label{節：発展・特殊型：解説}
\begin{enumerate}[label=\textbf{問\arabic*.},leftmargin=3.8em,format=\bookitemlink{practice-special}{solution}{problem}]
  % E1
  \item \textbf{(1)~整数部分を分ける.}
    \sectionref{節：ガウス記号を含む式}の整数性に着目する.
    実数$x$について,整数$\lfloor x\rfloor$を
    \[
      \lfloor x\rfloor=2q+r,\qquad r=0,1
    \]
    と表す.\bookterm{floor}{ガウス記号}の定義を使って$2$で割れば,
    \begin{align*}
      \frac x2&\ge q+\frac r2\ge q,\\
      \frac x2&<q+\frac{r+1}{2}\le q+1.
    \end{align*}
    したがって,
    \[
      \left\lfloor\frac{x}{2}\right\rfloor=q
      =\left\lfloor\frac{\lfloor x\rfloor}{2}\right\rfloor.
    \]
    $a_1=\lfloor2026\rfloor$であり,この恒等式を繰り返し使うと,
    \[
      a_n=\left\lfloor\frac{2026}{2^{n-1}}\right\rfloor
    \]
    が帰納法で得られる.

    \smallskip
    \textbf{(2)~1未満になる位置.}
    $2^{10}=1024\le2026<2048=2^{11}$だから,
    \[
      a_{11}=1,\qquad a_{12}=0.
    \]
    $n\le11$では切り捨てる前の値が$1$以上なので,
    初めて$0$となるのは$n=12$である.
    毎回の切り捨てを一つの\bookterm{floor}{ガウス記号}にまとめられる理由は,割る数が正の整数であることにある.

    \smallskip
    \textbf{(3)~整数を引いて同じ操作にする.}
    $c_n=b_n-6$とおく.整数をガウス記号の外へ出すと,
    \[
      c_{n+1}=\left\lfloor\frac{c_n}{2}\right\rfloor,
      \qquad c_1=M-6.
    \]
    (1)の恒等式は負の$x$にも成り立つので,
    \[
      c_n=\left\lfloor\frac{M-6}{2^{n-1}}\right\rfloor
    \]
    となる.十分大きな$n$では,中身が$M\ge6$なら$[0,1)$,$M<6$なら$(-1,0)$に入る.
    したがって$b_n$は,それぞれ$6$,$5$に有限回で到達する.
    実際,整数$x$が不動点となる条件は
    \[
      x\le\frac{x+6}{2}<x+1
      \iff 4<x\le6
    \]
    である.丸めを外した実数の漸化式の不動点は$6$だけだが,切り捨てると$5$も不動点となる.
    % E2
  \item \textbf{(1)~同じ値をまとめて足す.}
    \sectionref{節：同じ値を群にまとめる}の考え方を使う.
    \bookterm{difference}{階差}を足すと,
    \[
      a_n=\sum_{k=1}^{n}\left\lfloor\frac{k}{3}\right\rfloor.
    \]
    $k=1,2$では0で,$k=3j,3j+1,3j+2$では値が$j$になる.
    $n=3q+r$, $q\ge1$なら,値$1,\ldots,q-1$の群は三項ずつ,最後の値$q$の群は$r+1$項ある.
    よって,
    \begin{align*}
      a_n&=3\sum_{j=1}^{q-1}j+q(r+1)\\
      &=\frac{3q(q-1)}2+q(r+1).
    \end{align*}
    $q=0$でも実際の値と式は0である.
    最後の群は途中までなので,常に三項と数えてはいけない.

    \smallskip
    \textbf{(2)~和を個数として読む.}
    $x=k$のとき,指定範囲の整数$y$は$1$から$\lfloor\frac{k}{3}\rfloor$まである.
    $k=1,\ldots,n$について足せば,(1)の和に一致する.
    群ごとに値と項数を掛けて足す方法と,格子点の個数を数える方法は,同じ和の二つの見方である.
    % E3
  \item \textbf{(1)~余りに読み替える.}
    \sectionref{節：ガウス記号を含む式}では,\bookterm{floor}{ガウス記号}が整数問題の道具にもなる.
    整数$x$について\bookterm{floor}{ガウス記号}の定義から,
    \[
      0\le x-5\left\lfloor\frac{x}{5}\right\rfloor<5.
    \]
    整数なので値は$0,1,2,3,4$に限られる.
    初項の整数性と$x=2a_n+1$を使えば,帰納的に全項がこの範囲に入る.

    \smallskip
    \textbf{(2)~5つの値の移り方.}
    (1)より,次の表だけを調べれば十分である.
    \[
      \begin{array}{c|rrrrr}
        a_n&0&1&2&3&4\\ \hline
        a_{n+1}&1&3&0&2&4
      \end{array}
    \]
    初項$1$からは$1,3,2,0,1,\ldots$と進む.
    次は一意に決まるためこの四値を繰り返し,解答の\bookterm{general}{一般項}となる.
    四値が異なるので最小周期は$4$である.取り得る値を先に絞った.
    % E4
  \item \textbf{(1)~4項型から\bookterm{difference}{階差}へ.}
    \sectionref{節：隣接多項式}と同じく,\bookterm{difference}{階差}の組合せを見る.
    設問の$d_n$について,
    \begin{align*}
      d_{n+1}-d_n
      &=a_{n+3}-3a_{n+2}\\
      &\quad+3a_{n+1}-a_n\\
      &=6.
    \end{align*}
    $d_1=28-2\cdot9+2=12$だから,
    \[
      d_n=12+6(n-1)=6n+6.
    \]

    \smallskip
    \textbf{(2)~\bookterm{difference}{階差}を足して戻す.}
    $e_n=a_{n+1}-a_n$とおくと,
    \[
      e_{n+1}-e_n=d_n=6n+6,
    \]
    かつ$e_1=9-2=7$である.したがって,
    \begin{align*}
      e_n&=7+\sum_{k=1}^{n-1}(6k+6)\\
      &=3n^2+3n+1.
    \end{align*}
    もう一度足せば,
    \begin{align*}
      a_n&=2+\sum_{k=1}^{n-1}(3k^2+3k+1)\\
      &=n^3+1.
    \end{align*}
    ここでは$n=1$の空の和を$0$とする.
    初期値を$d_1,e_1,a_1$の順に使った.
    \bookterm{difference}{階差}で次数を下げ,和で元へ戻す.
    % E5
  \item \textbf{(1)~3項をまとめて等比型へ.}
    \sectionref{節：隣接多項式}の方法を使う.
    \bookterm{characteristic}{特性方程式}の左辺は
    \[
      x^3-6x^2+11x-6=(x-1)(x-2)(x-3)
    \]
    であり,$1,2$に対応する二次式から設問の$d_n$が作られる.
    実際に展開すると,
    \[
      d_{n+1}-3d_n
      =a_{n+3}-6a_{n+2}+11a_{n+1}-6a_n=0.
    \]
    $d_1=14-3\cdot6+2\cdot3=2$より,
    \[
      d_n=2\cdot3^{n-1}.
    \]
    初期値が0でないので,右辺を残して進める.

    \smallskip
    \textbf{(2)~右辺を作る\bookterm{sequence}{数列}を引く.}
    $3^{n-1}$に同じ組合せを施すと,
    \[
      3^{n+1}-3\cdot3^n+2\cdot3^{n-1}=2\cdot3^{n-1}.
    \]
    そこで$b_n=a_n-3^{n-1}$とおけば,
    \[
      b_{n+2}-3b_{n+1}+2b_n=0.
    \]
    特性方程式の解は$1,2$だから$b_n=A+B2^{n-1}$とおける.
    $b_1=2,b_2=3$から$A=B=1$となり,
    \[
      a_n=1+2^{n-1}+3^{n-1}
    \]
    を得る.4項間を3項間へ,さらに既習の解法へ順に帰着させた.
    % E6
  \item \textbf{(1)~特性方程式の解を角度で表す.}
    \sectionref{節：特性方程式の解が虚数の場合}を使う.
    \bookterm{characteristic}{特性方程式}$x^2-x+1=0$の根は,
    \[
      \frac{1\pm\sqrt3\,i}{2}
      =\cos\frac\pi3\pm i\sin\frac\pi3.
    \]
    絶対値が$1$なので,
    \[
      a_n=C\cos\frac{(n-1)\pi}{3}
      +D\sin\frac{(n-1)\pi}{3}
    \]
    とおける.$a_1=0$より$C=0$であり,
    $a_2=1$より$D=\frac2{\sqrt3}$である.
    得られた漸化式は$\sin$の\bookterm{addition}{加法定理}からも確認できる.

    \smallskip
    \textbf{(2)~最小周期を確かめる.}
    \bookterm{general}{一般項}より$a_{n+6}=a_n$であり,最初の6項は
    \[
      0,1,1,0,-1,-1
    \]
    である.周期が$1,2,4,5$なら$a_{1+p}=a_1$が成り立たない.
    周期が$3$なら$a_5=a_2$が必要だが,
    $-1\ne1$である.
    したがって周期の最小値は$6$である.

    \smallskip
    \noindent \textbf{【別解】}\\
    \bookterm{recurrence}{漸化式}を一つ進めると,
    \[
      a_{n+3}=a_{n+2}-a_{n+1}=-a_n.
    \]
    さらに3項進めれば$a_{n+6}=a_n$となる.
    特性方程式の解の角度と,直接の式変形による周期の証明が対応している.
    % E7
  \item \textbf{振幅と\bookterm{limit}{極限}.}
    \sectionref{節：特性方程式の解が虚数の場合}では,特性方程式の解の絶対値も見る.
    \bookterm{characteristic}{特性方程式}$x^2-x+\frac12=0$の根は,
    \[
      \frac{1\pm i}{2}
      =\frac1{\sqrt2}
      \left(\cos\frac\pi4\pm i\sin\frac\pi4\right).
    \]
    したがって$t_n=\frac{(n-1)\pi}{4}$として,
    \[
      a_n=\left(\frac1{\sqrt2}\right)^{n-1}
      (C\cos t_n+D\sin t_n)
    \]
    とおける.$a_1=1$から$C=1$であり,
    $a_2=0$から$D=-1$である.
    よって解答の\bookterm{general}{一般項}を得る.\bookterm{trigonometric}{三角関数}の絶対値は$1$以下なので,
    \[
      |a_n|\le2\left(\frac1{\sqrt2}\right)^{n-1}\longrightarrow0.
    \]
    したがって$a_n\to0$となる.
    角度は一定ずつ進み,振幅は\bookterm{geometric}{等比数列}として縮んでいる.
    % E8
  \item \textbf{(1)~正値の確認と奇偶の分割.}
    \sectionref{節：非隣接2項型}と同じく奇偶に分ける.
    $a_1,a_2>0$で,正の項から次も正になるので両方で帰納法を使えば全項が正である.
    逆数を取ると,
    \[
      \frac1{a_{n+2}}=\frac1{a_n}+1.
    \]
    $b_n=\frac1{a_n}$とおくと,
    \[
      b_1=1,\quad b_2=\frac12,\quad b_{n+2}=b_n+1.
    \]
    したがって,
    \[
      b_{2k-1}=k,\qquad b_{2k}=k-\frac12.
    \]
    逆数を取って解答の\bookterm{general}{一般項}を得る.

    \smallskip
    \textbf{(2)~部分列から\bookterm{limit}{極限}へ.}
    奇数番目と偶数番目のどちらも$0$へ近づく.
    全体としても,$n\ge2$では\bookterm{general}{一般項}から
    \[
      0<a_n\le\frac2{n-1}
    \]
    となるので,$a_n\to0$である.
    添字を分けたあとには,各部分列の答えを元の添字へ戻して確認する.
    % E9
  \item \textbf{(1)~値をずらす.}
    \bookterm{square}{平方完成}して$b_n=a_n-1$とおくと,
    \begin{align*}
      b_{n+1}&=a_n^2-2a_n-1\\
      &=(a_n-1)^2-2=b_n^2-2.
    \end{align*}
    初項は$b_1=\frac52$である.

    \smallskip
    \textbf{(2)~対称式の指数を変える.}
    \sectionref{節：対称式を使う}と同じ形が現れた.
    $\frac52=2+2^{-1}$であるから,
    \[
      B_n=2^{\,2^{n-1}}+2^{-2^{n-1}}
    \]
    を候補とする.設問の恒等式に$u=2^{\,2^{n-1}}$を入れると,
    \[
      B_{n+1}=B_n^2-2.
    \]
    $B_1=b_1$でもあるので,帰納法により$b_n=B_n$となる.
    よって$a_n=1+B_n$である.
    \bookterm{square}{平方完成}の後,逆数となる二量の和として恒等式を使った.
    % E10
  \item \textbf{(1)~交差項の打ち消し.}
    \sectionref{節：対称式を使う}の三次の恒等式を使う.
    展開すると,
    \[
      (u+u^{-1})^3=u^3+3u+3u^{-1}+u^{-3}
    \]
    であるから,$3(u+u^{-1})$を引けばよい.

    \smallskip
    \textbf{(2)~指数を3倍する.}
    $a_1=2+2^{-1}$に合う候補を
    \[
      A_n=2^{\,3^{n-1}}+2^{-3^{n-1}}
    \]
    とする.(1)に$u=2^{\,3^{n-1}}$を代入すると,
    \[
      A_{n+1}=A_n^3-3A_n
    \]
    となる.初項と\bookterm{recurrence}{漸化式}が一致するので,帰納法より$a_n=A_n$である.
    $-3a_n$の項が交差項を正確に打ち消すため,同じ対称式の形を保って次へ進める.
    % E11
  \item \textbf{(1)~倍角から候補を作る.}
    \sectionref{節：三角関数型}の余弦の倍角公式を使う.
    $\theta_n=2^{n-1}\frac{2\pi}{7}$とおき,
    $A_n=\cos\theta_n$とすれば,
    \[
      A_{n+1}=\cos(2\theta_n)=2A_n^2-1.
    \]
    初項も一致するので,帰納法により$a_n=A_n$である.
    角度から解を作っており,余弦の値から角度を一意に決めたのではない.

    \smallskip
    \textbf{(2)~一周の差を比べる.}
    3回進めると,
    \[
      \theta_{n+3}-\theta_n=2^n\pi
      =2\pi\cdot2^{n-1}.
    \]
    これは$2\pi$の整数倍だから,$a_{n+3}=a_n$である.
    $\cos\frac{8\pi}{7}=\cos\left(2\pi-\frac{8\pi}{7}\right)=\cos\frac{6\pi}{7}$に注意すると,最初の三つの値は,
    \[
      \cos\frac{2\pi}{7},\quad
      \cos\frac{4\pi}{7},\quad
      \cos\frac{6\pi}{7}
    \]
    となる.$\cos$は$0\le\theta\le\pi$で狭義に減少するので,
    三つの値は異なり,周期の最小値は$3$である.
    % E12
  \item \textbf{(1)~二つの分母を確認する.}
    \sectionref{節：三角関数型}の正接の倍角公式から,
    \[
      \theta_n=2^{n-1}\frac\pi5,\qquad A_n=\tan\theta_n
    \]
    を候補とする.$\cos\theta_n=0$となるには,
    ある整数$j$について
    \[
      2^{n-1}\frac\pi5=\frac\pi2+j\pi,
      \quad\text{すなわち}\quad 2^n=5(2j+1)
    \]
    が必要である.左辺は偶数,右辺は奇数なので,これは起こらない.
    よって$A_n$はすべての$n\ge1$で定義できる.
    さらに,
    \[
      1-A_n^2=\frac{\cos(2\theta_n)}{\cos^2\theta_n}
      =\frac{\cos\theta_{n+1}}{\cos^2\theta_n}\ne0.
    \]
    したがって,倍角公式を各段階で使うことができ,
    \[
      A_{n+1}=\frac{2A_n}{1-A_n^2}
    \]
    となる.初項も一致するから,帰納法で$a_n=A_n$を得る.

    \smallskip
    \textbf{(2)~正接の周期と角度の差.}
    4回進めると,
    \[
      \theta_{n+4}-\theta_n=3\cdot2^{n-1}\pi.
    \]
    $\tan$の周期は$\pi$なので,$a_{n+4}=a_n$である.
    最初の四つの値は,
    \[
      \tan\frac\pi5,\quad \tan\frac{2\pi}{5},\quad
      -\tan\frac\pi5,\quad -\tan\frac{2\pi}{5}
    \]
    である.$0<\frac\pi5<\frac{2\pi}{5}<\frac\pi2$で正接は狭義に増加するから,
    この四つは異なる.よって周期の最小値は$4$である.
    % E13
  \item \textbf{(1)~正値を確かめる.}
    $a_1=1>0$であり,$a_n>0$なら分子・分母がともに正である.
    よって帰納法によりすべての項が正であり,$a_n+\sqrt3$も$0$にならない.

    \smallskip
    \textbf{(2)~二つの平方の比.}
    \sectionref{節：ニュートン法}の$\sqrt2$を$\sqrt3$に置き換えた形である.
    \[
      a_{n+1}-\sqrt3=\frac{(a_n-\sqrt3)^2}{2a_n},
    \]
    \[
      a_{n+1}+\sqrt3=\frac{(a_n+\sqrt3)^2}{2a_n}.
    \]
    共通の分母を消すために比を取ると,
    \[
      b_{n+1}=b_n^2.
    \]
    初項は
    \[
      b_1=\frac{1-\sqrt3}{1+\sqrt3}=\sqrt3-2=q.
    \]
    $q<0$なので最初から実数の\bookterm{logarithm}{対数}は使えない.二乗の繰り返しから帰納的に
    \[
      b_n=q^{\,2^{n-1}}
    \]
    が得られる.$b_n$の定義を$a_n$について解けば,
    \[
      a_n=\sqrt3\,\frac{1+b_n}{1-b_n}.
    \]
    $-1<q<0$なので$|b_n|<1$であり,
    分母$1-b_n$は常に正である.

    \smallskip
    \textbf{(3)~第2項以降の減少.}
    $n\ge2$では$b_n>0$なので,\bookterm{general}{一般項}から$a_n>\sqrt3$である.
    また,
    \[
      a_{n+1}-a_n=\frac{3-a_n^2}{2a_n}<0
      \quad(n\ge2).
    \]
    $|q|<1$から$b_n\to0$であり,$a_n\to\sqrt3$となる.
    初項は収束先の下にあるが,第2項で上へ移ってから減少する.
    誤差を二乗する式も,この移動と速い収束を表す.
    % E14
  \item \textbf{(1)~\bookterm{fixed}{不動点}と定\bookterm{sequence}{数列}.}
    \sectionref{節：初項による違い}の考え方を使う.
    定\bookterm{sequence}{数列}には$c=c^2-1$が必要だから,
    \[
      c^2-c-1=0,\qquad c=\frac{1\pm\sqrt5}{2}.
    \]
    これらの初項なら毎回同じ値になるので,十分条件でもある.

    \smallskip
    \textbf{(2)~二回で戻る初項.}
    $a_3=a_1$から,
    \begin{align*}
      (c^2-1)^2-1&=c,\\
      c^4-2c^2-c&=0,\\
      c(c+1)(c^2-c-1)&=0.
    \end{align*}
    (1)の定\bookterm{sequence}{数列}を除くと$c=0,-1$だけが残る.
    実際,項は$0\mapsto-1\mapsto0$と移るので,
    どちらの初項もこの二つの値を交互に取る.
    その\bookterm{general}{一般項}は
    \[
      a_n=-\frac12+\left(c+\frac12\right)(-1)^{n-1}
      \quad(c=0,-1)
    \]
    と表せる.二つの値が異なるから周期の最小値は$2$である.
    ここでの\bookterm{general}{一般項}は,調べた初項に限って使える.
    % E15
  \item \textbf{(1)~初項による振る舞い.}
    代入すると解答の表を得る.$c=-1,2$は初めから定\bookterm{sequence}{数列},$c=-2,0$は数項後に$2$で一定になる.

    \smallskip
    \textbf{(2)~初項に合う表示.}
    \sectionref{節：初項による違い}から\sectionref{節：三角関数型}へつながる例である.
    $-2\le c\le2$なら$c=2\cos\theta$
    （$0\le\theta\le\pi$）と表せる.
    $A_n=2\cos(2^{n-1}\theta)$とおくと,
    \begin{align*}
      A_n^2-2
      &=4\cos^2(2^{n-1}\theta)-2\\
      &=2\cos(2^n\theta)=A_{n+1}.
    \end{align*}
    初項も一致するので,帰納法により$a_n=A_n$である.
    $|\cos x|\le1$から,すべての項で$|a_n|\le2$となる.

    \smallskip
    \textbf{(3)~2からの距離を評価する.}
    $|c|>2$なら$a_2=c^2-2>2$である.
    $a_n>2$のとき,
    \[
      a_{n+1}-2=(a_n-2)(a_n+2)>4(a_n-2)>0.
    \]
    よって帰納法で$a_n>2$が$n\ge2$のすべてで続く.
    この不等式を繰り返すと,
    \[
      a_n-2\ge4^{n-2}(a_2-2)
      =4^{n-2}(c^2-4).
    \]
    右辺は$+\infty$へ発散するので,$a_n\to+\infty$である.
    したがって有界となる初項はちょうど$-2\le c\le2$である.
    \bookterm{general}{一般項}が表せる範囲と,不等式で振る舞いを調べる範囲を,初項に応じて選んだ.
\end{enumerate}
